\begin{pastexam}[label={px:ust-115-ee-ds-13}]{臺聯大電機類 115 學年度 第 13 題}
（單選；2 分）

Skip List is a probabilistic data structure based on linked lists. What is its main purpose in
addressing the shortcomings of standard linked lists?
\begin{enumerate}[label=(\Alph*)]
\item Solving the problem of slow insertion
\item Solving the problem of not being able to perform fast searches
\item Solving the problem of excessive memory usage
\item Solving the problem of being unable to traverse in reverse
\end{enumerate}
\end{pastexam}

\begin{pastsoln}
答案是 (B)。

排序好的串列不能隨機存取，找一個鍵值只能從頭一個一個走，要 $\Theta(n)$，不能像陣列那樣二分搜尋。
跳躍串列在原來的串列上面加幾層「快速道路」：每個節點隨機決定高度，第 $i$ 層大約每 $2^i$ 個節點
留一個。 搜尋從最高層開始往右走，走過頭就往下一層，期望時間是 $O(\log n)$（資結書第~10.7~節、
定理~10.7）。 它的目的是讓串列也能快速搜尋。

逐項檢查：
\begin{itemize}
\item (A) 串列在已知位置插入本來就是 $\Theta(1)$；慢的是「找到插入位置」，那正是搜尋的問題。
跳躍串列的插入也要先搜尋。 主要目的不是插入。 錯。
\item (B) 對。
\item (C) 跳躍串列多了每層的指標，期望大約 $2n$ 個指標，記憶體反而比單一串列多。 錯。
\item (D) 往回走要靠雙向串列，跳躍串列的指標都是往右的。 錯。
\end{itemize}
\end{pastsoln}
